\section{Prompt Injection Attacks}
\label{sec:prompt-injection}

Prompt injection occurs when attacker-controlled text is interpreted as an
instruction that overrides or redirects the trusted task
\parencite{perez2022ignore}. This does not necessarily require a conventional
software defect. The failure arises because an instruction-following model can
assign authority to text from an unintended source. The surrounding application
turns that behavioural failure into a security problem when it exposes protected
data or permits consequential actions.

\subsection{Direct and Indirect Prompt Injection}

The literature distinguishes two principal delivery mechanisms
\parencite{greshake2023indirect}.

In a \emph{direct prompt injection}, a malicious user submits the payload through
the target application's ordinary user-input channel. The attacker may ask the
model to ignore the established task, imitate a higher-priority instruction, or
begin a new task
\parencite{perez2022ignore}. Figure~\ref{fig:direct-prompt-injection} shows both
the trusted instruction and the competing attacker-controlled input.

\begin{figure}[htbp]
    \centering
    \begin{tikzpicture}[
        box/.style={draw, rounded corners, align=center, text width=2.5cm,
                    minimum height=0.9cm},
        untrusted/.style={box, dashed},
        flow/.style={-{Latex[length=2.2mm]}, thick}
    ]
        \node[box] (trusted) at (1.4,1.3)
            {Trusted application\\instruction};
        \node[untrusted] (attacker) at (1.4,-1.3)
            {Malicious user's\\direct input};
        \node[box] (application) at (5.0,0)
            {Application\\prompt construction};
        \node[box] (llm) at (8.7,0)
            {Large Language\\Model};
        \node[box] (output) at (12.4,0)
            {Model\\output};

        \draw[flow] (trusted) -- (application);
        \draw[flow] (attacker) -- (application);
        \draw[flow] (application) -- (llm);
        \draw[flow] (llm) -- (output);
    \end{tikzpicture}
    \caption[Direct prompt injection]{Direct prompt injection through the
    application's user-input channel. The malicious input is combined with the
    trusted application instruction during prompt construction. Original
    schematic based on \textcite{perez2022ignore}.}
    \label{fig:direct-prompt-injection}
\end{figure}

In an \emph{indirect prompt injection}, the payload is delivered through content
that the application retrieves or receives during normal operation, such as a
web page, document, email, application programming interface (API) response, or
tool output. The attacker need not submit the completed payload directly to the
target application. This makes indirect injection comparatively covert, but not
inherently more damaging: the impact still depends on the data and authority
available to the application
\parencites{greshake2023indirect}{debenedetti2024agentdojo}.

Retrieval-augmented generation (RAG) systems and tool-using agents are exposed
because they routinely place external content in the model context. As
Figure~\ref{fig:indirect-prompt-injection} illustrates, trusted and untrusted
inputs converge at the application even though they originate from different
principals.

\begin{figure}[htbp]
    \centering
    \begin{tikzpicture}[
        box/.style={draw, rounded corners, align=center, text width=2.3cm,
                    minimum height=0.9cm},
        untrusted/.style={box, dashed},
        flow/.style={-{Latex[length=2.2mm]}, thick}
    ]
        \node[box] (trusted) at (1.3,1.4)
            {Trusted task\\instruction};
        \node[untrusted] (content) at (1.3,-1.4)
            {External content\\with injected\\instruction};
        \node[box] (retrieval) at (4.4,-1.4)
            {Retrieval system,\\API, or tool};
        \node[box] (application) at (7.5,0)
            {Application\\context assembly};
        \node[box] (llm) at (10.5,0)
            {Large Language\\Model};
        \node[box] (effect) at (13.5,0)
            {Output or\\tool request};

        \draw[flow] (trusted) -- (application);
        \draw[flow] (content) -- (retrieval);
        \draw[flow] (retrieval) -- (application);
        \draw[flow] (application) -- (llm);
        \draw[flow] (llm) -- (effect);
    \end{tikzpicture}
    \caption[Indirect prompt injection]{Indirect prompt injection through
    attacker-controlled external content retrieved at run time. The content is
    incorporated into the application context together with the trusted task
    instruction. Original schematic based on
    \textcite{greshake2023indirect}.}
    \label{fig:indirect-prompt-injection}
\end{figure}

Embedding an instruction in content processed at run time is distinct from
training-data poisoning. It is also narrower than poisoning a retrieval index or
manipulating a retriever's ranking. Those attacks may determine which content is
selected, whereas indirect prompt injection concerns how malicious content
influences the model after it enters the inference context.

\subsection{Why Prompt Injection Works and Its Consequences}

Prompt and role formats provide useful signals, but current defences do not
justify assuming that those signals form a reliably enforced provenance or
authority boundary. A model can therefore respond to instructional language in
an untrusted channel even when the application intended it to be data
\parencites{liu2023promptinjection}{wallace2024instruction}. Instruction
following is both the capability that enables useful task execution and the
behaviour an injected payload attempts to redirect.

The consequences depend on the surrounding system. The Prompt Injection
Benchmark for Foundation Model Integrated Systems (FSPIB) distinguishes
\emph{information-level} threats, including information leakage, goal hijacking,
and response refusal, from \emph{action-level} threats such as adversarial
action and parameter manipulation
\parencite{fspib2025}. The distinction separates control over generated text
from control that propagates into a connected application or tool.

Prompt injection can therefore bridge natural-language manipulation and more
traditional attack surfaces
\parencite{pedro2025p2sql}. In a 2024 researcher-reported Slack AI
demonstration, a malicious instruction in a public channel caused retrieved
private-channel information to be encoded in a generated link. Exfiltration
occurred only if the victim clicked that attacker-controlled link, so the case
demonstrates an interaction-dependent attack chain rather than an independently
confirmed no-click breach
\parencite{promptarmor2024slack}.

\subsection{Attack Taxonomy}

Prompt-injection categories are clearest when treated as orthogonal axes rather
than mutually exclusive labels. Table~\ref{tab:attack-taxonomy} separates these
dimensions.

\begin{table}[htbp]
    \centering
    \small
    \caption{Orthogonal dimensions of a prompt injection attack.}
    \label{tab:attack-taxonomy}
    \begin{tabularx}{\linewidth}{@{}p{3.0cm}YY@{}}
        \toprule
        \textbf{Dimension} & \textbf{Representative categories} &
        \textbf{Question answered} \\
        \midrule
        Delivery or source & Direct user input; indirect document, web, email,
        API, tool, or memory content & How does the payload enter the target
        context? \\
        Construction method & Handcrafted heuristic; task-specific
        optimisation; universal or transferable suffix & How is the payload
        produced? \\
        Capability and knowledge & Query access; weights or gradients; context
        knowledge; defence knowledge & What access and information are used
        during construction? \\
        Objective or impact & Instruction-integrity violation; disclosure;
        unauthorised tool action; availability loss & What security property is
        the attacker trying to violate? \\
        \bottomrule
    \end{tabularx}
\end{table}

Heuristic attacks include context-ignoring and fake-completion prompts
\parencite{jia2025critical}. A context-ignoring attack asks the model to discard
the established task. A completion attack imitates the end of the legitimate
exchange before introducing a new instruction. Both are manually constructed,
but either may be delivered directly or indirectly.

Optimisation-based methods search for token sequences that increase the
likelihood of an attacker-selected output. Greedy Coordinate Gradient (GCG) is
a prominent example. It was introduced as an adversarial-suffix method for
jailbreaking safety-aligned models, not specifically for prompt injection
\parencite{zou2023universal}. StruQ and SecAlign adapt its objective to make a
model follow an injected task instead of the trusted task
\parencites{chen2025struq}{chen2025secalign}. GCG uses model gradients while
constructing the suffix, although the resulting payload may later be delivered
through an ordinary indirect channel.

Jailbreaking and prompt injection should consequently remain distinct. A
jailbreak attempts to override a provider's safety policy, whereas prompt
injection attempts to redirect an application- or user-defined trusted task
through attacker-controlled data. The same optimisation technique can be used
for both objectives without making the threat models identical. One attack can,
for example, be indirectly delivered, optimisation-generated with white-box
access, constructed without knowledge of the victim's context, and aimed at data
exfiltration.
